\subsection{李群的自同构群（实或复情形）}

本号中假设 K = R 或 C。

引理 4. 设 \(\mathbf{H}\) 是有限维单连通李群。

\begin{enumerate}
\def\labelenumi{(\roman{enumi})}
\item
  对所有 \(u \in \operatorname{Aut} \mathbf{L}(\mathbf{H})\)，令 \(\theta(u)\) 为 \(\mathbf{H}\) 的唯一自同构，使得 \(\mathbf{L}(\theta(u)) = u\)。则映射 \((u, g) \mapsto \theta(u)g\) 从 \((\operatorname{Aut} \mathbf{L}(\mathbf{H})) \times \mathbf{H}\) 到 \(\mathbf{H}\) 是解析的。
\item
  设 \(\mathbf{N}\) 是 \(\mathbf{H}\) 的李子群，\(\operatorname{Aut}(\mathbf{H},\mathbf{N})\) 是由 \(v\in \operatorname{Aut}\mathbf{H}\) 组成的集合，其中 \(v(\mathbf{N}) = \mathbf{N}\)。则 \(\theta^{-1}(\operatorname {Aut}(\mathbf{H},\mathbf{N}))\) 是 \(\operatorname {Aut}\mathbf{L}(\mathbf{H})\) 的李子群。
\item
  假设 N 是离散正规子群，从而 G = H/N 的李代数可与 L(H) 识别。对所有 \(w \in Aut G\)，令 \(\eta(w)\) 为 H 的唯一自同构，使得 \(\mathrm{L}(\eta(w)) = \mathrm{L}(w)\)。则映射 \(\eta\) 是从群 Aut G 到群 \(\mathrm{Aut}(\mathrm{H}, \mathrm{N})\) 的同构。
\end{enumerate}

要证明 (i)，根据第 1 号引理 3，只需验证映射 \((u,g)\mapsto\theta(u)g\) 在 \((\mathrm{Id}_{\mathrm{L(H)}},e)\) 的一个邻域内是解析的。存在 L(H) 中 0 的一个开邻域 B，使得 \(\psi = \exp_{\mathbf{H}}|\mathbf{B}\) 是从 B 到 H 中 \(e\) 的一个开邻域的解析同构。存在 Aut L(H) 中 \(\mathrm{Id}_{\mathrm{L(H)}}\) 的一个开邻域 U 和 L(H) 中 0 的一个开邻域 B'，使得 \(\mathrm{U(B')} \subset \mathrm{B}\)。于是映射 \((u, g) \mapsto \theta(u)g\) 从 \(\mathrm{U} \times \psi(\mathrm{B'})\) 到 H 是由以下映射复合而成：

\[
\begin{array}{l l l l} \text {the mapping} (u, g) \mapsto \langle u, \psi^ {- 1} (g) \rangle & \text {of} & \mathrm{U} \times \psi (\mathrm{B} ^ {\prime}) & \text {into} \quad \mathrm{U} \times \mathrm{B} ^ {\prime}; \\ \text {the mapping} (u, x) \mapsto u (x) & \text {of} & \mathrm{U} \times \mathrm{B} ^ {\prime} & \text {into} \quad \mathrm{B}; \\ \text {the mapping} y \mapsto \psi (y) & \text {of} & \mathrm{B} & \text {into} \quad \mathrm{G}. \end{array}
\]

所以这个映射是解析的。

令 \(p\) 为从 \(H\) 到齐性空间 \(H/N\) 的典范映射。则 \(\theta^{-1}(\text{Aut}(H, N))\) 是满足以下条件的 \(u \in \text{Aut } L(H)\) 的集合：

\[
p (\theta (u) g) = p (e), \quad p (\theta (u ^ {- 1}) g) = p (e)
\]

对所有 \(g \in \mathbb{N}\) 均成立。由第 8 节第 2 号定理 2 及其推论 2，这就证明了 (ii)。

假设 N 是离散正规子群。令 \(w \in \operatorname{Aut} G\)。于是

\[
\mathrm{L} (p \circ \eta (w)) = \mathrm{L} (\eta (w)) = \mathrm{L} (w) = \mathrm{L} (w \circ p)
\]

所以 \(p \circ \eta(w) = w \circ p\)，从而 \(\eta(w) \in \text{Aut(H, N)}\)。显然，从 Aut G 到 Aut(H, N) 的映射 \(\eta\) 是单同态。这个同态是满射的，因为 \(p: H \to G\) 是浸没。

设 G 是局部紧群，\(\Gamma\) 是 G 的自同构群。回忆一下，\(T_{\beta}\) 上已经定义了一个拓扑 \(\Gamma\)（《一般拓扑学》第十章第 3 节第 5 号）。它是使映射 \(v \mapsto v\) 和 \(v \mapsto v^{-1}\) 连续的最粗拓扑，其中这些映射从 \(\Gamma\) 到 \(\mathcal{C}_{\varepsilon}(G; G)\)。拓扑 \(T_{\beta}\) 与 \(\Gamma\) 上的群结构相容（同上）。对 G 的每个紧子集 L 和 G 中 \(e_{\mathrm{G}}\) 的每个邻域 U，令 N(L, U) 为满足 \(\phi \in \Gamma\)、\(\phi(g) \in gU\)、\(\phi^{-1}(g) \in gU\) 对所有 \(g \in L\) 成立的集合；则 N(L, U) 构成 \(e_{\Gamma}\) 的邻域基。若 G 由一个紧子集 C 生成，则 \(T_{\beta}\) 也是使映射 \(v \mapsto v[C\) 和 \(v \mapsto v^{-1}[C]\) 从 \(\Gamma\) 到 \(\mathcal{C}_{u}(C; G)\) 连续的最粗拓扑（因为 G 的每个紧子集都包含于 \((\mathrm{C} \cup \mathrm{C}^{-1})^{\mathrm{n}}\)（当 n 足够大时））。若 K 是局部紧的，V 是 K 上的有限维向量空间，则拓扑 \(T_{\beta}\) 正是通常拓扑。

定理 1. 设 G 是有限维李群，\(G_{0}\) 是其单位元分支。假设 G 由 \(G_{0}\) 和有限个元素生成。

\begin{enumerate}
\def\labelenumi{(\roman{enumi})}
\tightlist
\item
  Aut G 上存在唯一的解析流形结构，满足以下条件：
\end{enumerate}

(AUT) 对每个解析流形 M 及每个从 M 到 Aut G 的映射 f，f 是解析的，当且仅当从 M × G 到 G 的映射 \((m, g) \mapsto f(m)g\) 是解析的。

以下均假设 Aut G 具有这一结构。

\begin{enumerate}
\def\labelenumi{(\roman{enumi})}
\setcounter{enumi}{1}
\item
  Aut G 是有限维李群。
\item
  从 Aut G 到 Aut L(G) 的态射 \(\phi: u \mapsto L(u)\) 是解析的。
\item
  若 G 是连通的，则 \(\phi\) 是从李群 Aut G 到 Aut L(G) 的一个李子群的同构；若 G 是单连通的，这个李子群等于 Aut L(G)。
\item
  令 \(\mathfrak{a}\) 为 \(\mathbf{G}\) 的无穷小自同构的集合。则 \(\mathfrak{a}\) 是一个向量场李代数，且与映射 \((u, g) \mapsto u(g)\) 相应的无穷小运算律（该映射从 (Aut G) \(\times\) G 到 G）是从 L(Aut G) 到 \(\mathfrak{a}\) 的同构。
\item
  李群 Aut G 的拓扑就是拓扑 \(\mathcal{T}_{\beta}\)。
\end{enumerate}

\begin{enumerate}
\def\labelenumi{(\alph{enumi})}
\item
  (i) 中所考虑的解析结构的唯一性是显然的。
\item
  假设 G 是连通的。令 H 为 G 的万有覆叠空间，p 为从 H 到 G 的典范态射，且 N = Ker p。~我们引入引理 4 中的记号 \(\theta\)、\(\eta\) 和 \(\operatorname{Aut}(\mathrm{H},\mathrm{N})\)。借助 \(\theta\)，将 Aut L(G) 的李群结构传递给 Aut H。于是 Aut H 成为有限维李群，\(\operatorname{Aut}(\mathrm{H},\mathrm{N})\) 成为 Aut H 的李子群（引理 4 (ii)）。将 \(\operatorname{Aut}(\mathrm{H},\mathrm{N})\) 的李群结构传递给 Aut G，所用映射为 \(\eta^{-1}\)。于是 Aut G 成为有限维李群。定理的性质 (ii)、(iii) 和 (iv) 得到满足，并且映射 \((u,g)\mapsto u(g)\) 从 \((\operatorname{Aut}\mathbf{G})\times\mathbf{G}\) 到 G 是解析的（引理 4 (i)）。设 M 是解析流形，f 是从 M 到 Aut G 的映射，\(\phi\) 是从 M × G 到 G 的映射 \((m,g)\mapsto f(m)g\)。显然，若 f 是解析的，则 \(\phi\) 是解析的。假设 \(\phi\) 是解析的。于是 T \(\phi\) : TM × TG → TG 是解析的；它在 M × L(G) 上的限制，即映射 \((m,x)\mapsto\operatorname{L}(f(m))x\)，因而是解析的；由于 L(G) 是有限维的，可得映射 \(m\mapsto\operatorname{L}(f(m))\) 是解析的，从而 f 是解析的。因此 (i) 成立。
\end{enumerate}

给 \(\mathrm{L}(\mathrm{G})\) 赋予一个范数。对所有 \(\lambda > 0\)，令 \(\mathrm{B}_{\lambda}\) 为以 0 为中心、半径为 \(\lambda\) 的 \(\mathrm{L}(\mathrm{G})\) 中的开球。选择 \(\lambda > 0\) 充分小，使得 \(\psi = \exp_{\mathrm{G}}|\mathrm{B}_{\lambda}\) 是从解析流形 \(\mathrm{B}_{\lambda}\) 到 \(\psi(\mathrm{B}_{\lambda})\) 的 \(G\) 的开子流形的同构。令 \(\Phi\) 是 Aut \(G\) 上的滤子。对于 \(\Phi\)，它收敛到 \(\mathrm{Id}_{\mathrm{G}}\)，且是在 Aut \(G\) 中的收敛，充要条件是 \(\mathrm{L}(\Phi)\) 收敛到 \(\mathrm{Id}_{\mathrm{L}(\mathrm{G})}\)，且是在 Aut \(L(\mathrm{G})\) 中的收敛，从而 \(L(\Phi)|B_{\lambda/2}\) 和 \(L(\Phi)^{-1}\mathrm{B}_{\lambda/2}\) 一致收敛到 \(\mathrm{Id}_{\mathrm{B}_{\lambda/2}}\)。这一条件蕴含 \(\Phi|\psi(\mathrm{B}_{\lambda/2})\) 和 \(\Phi^{-1}|\psi(\mathrm{B}_{\lambda/2})\) 一致收敛到 \(\mathrm{Id}_{\psi(\mathrm{B}_{\lambda/2})}\)。反之，假设 \(\Phi|\psi(\mathrm{B}_{\lambda/2})\) 一致收敛到 \(\mathrm{Id}_{\psi(\mathrm{B}_{\lambda/2})}\)。存在 \(M \in \Phi\)，使得若 \(u \in M\)，则 \(u(\psi(\mathrm{B}_{\lambda/2})) \subset \psi(\mathrm{B}_{2\lambda/3})\)；于是 \(L(u)(\mathrm{B}_{\lambda/2})\) 是 \(L(\mathrm{G})\) 的连通子集，其在 \(\exp_{\mathrm{G}}\) 下的像包含于 \(\psi(\mathrm{B}_{2\lambda/3})\) 中，因此 \(L(u)(\mathrm{B}_{\lambda/2})\) 不与 \(B_{\lambda} - B_{2\lambda/3}\) 相交，从而 \(L(u)(\mathrm{B}_{\lambda/2}) \subset B_{\lambda}\)；于是 \(\Phi|\psi(\mathrm{B}_{\lambda/2})\) 一致收敛到 \(\mathrm{Id}_{\psi(\mathrm{B}_{\lambda/2})}\) 的假设蕴含 \(L(\Phi)|B_{\lambda/2}\) 一致收敛到 \(\mathrm{Id}_{\mathrm{B}_{\lambda/2}}\)。于是得到：

\((\Phi\) 在 Aut G 中收敛到 \(\mathrm{Id}_{\mathrm{G}}\)) \(\Leftrightarrow\)（\(\Phi\) 在 Aut G 中收敛到 \(\mathrm{Id}_{\mathrm{G}}\)，在 \(\mathcal{T}_{\beta}\) 下）。

这就证明了 (vi)。

令 D 为 \(\mathrm{Aut(G)}\) 在 G 上的左作用律所对应的无穷小运算律。由第 1 号命题 1 和 2，\(\mathrm{D(L(AutG))} = a\)。因此 \(a\) 是一个向量场李代数，且 D 是从 \(\mathrm{L(AutG)}\) 到 \(a\) 的态射。设 \(x_{1}\) 和 \(x_{2}\) 是 \(\mathrm{L(AutG)}\) 中满足 \(\mathrm{D}(x_1) = \mathrm{D}(x_2)\) 的元素。则 K 在 G 上的两个运算律 \((\lambda, g) \mapsto (\exp \lambda x_1)g\) 和 \((\lambda, g) \mapsto (\exp \lambda x_2)g\) 具有相同的相应无穷小运算律；因此，当 \(|\lambda|\) 充分小时，\(\exp \lambda x_1\) 和 \(\exp \lambda x_2\) 在 \(e\) 的一个邻域内重合（第 4 节第 7 号定理 6），从而 \(\exp \lambda x_1 = \exp \lambda x_2\)。可得 \(x_1 = x_2\)，所以 D 是从 \(\mathrm{L(AutG)}\) 到 \(a\) 的同构。

因此，定理在 G 连通时已经完全证明。

\begin{enumerate}
\def\labelenumi{(\alph{enumi})}
\setcounter{enumi}{2}
\tightlist
\item
  转到一般情形。根据假设，G 由 \(G_{0}\) 和有限个元素 \(x_{1}, x_{2}, \ldots, x_{n}\) 生成。每个 \(u \in AutG\) 都保持 \(G_{0}\) 稳定。令 \(Aut_{1}G\) 为这样的 \(u \in AutG\) 的集合：它们在取商后诱导出 \(G/G_{0}\) 的恒等自同构。这是 Aut G 的正规子群。由证明的 (b) 部分，Aut \(AutG_{0}\) 具有典范的李群结构，并且映射 \((g_{1}, g_{2}, \ldots, g_{n}, u) \mapsto (ug_{1}, ug_{2}, \ldots, ug_{n})\) 从 \(G_{0}^{n} \times AutG_{0}\) 到 \(G_{0}^{n}\) 是解析的。令 P 为 \(AutG_{0}\) 与 \(G_{0}^{n}\) 相应的半直积；它是有限维李群（第 \(\S\) 1 节第 4 号命题 7）。
\end{enumerate}

若 \(w\in \mathrm{Aut}_1\mathrm{G}\)，记

\[
\begin{array}{r l} w _ {0} & = w | \mathrm{G} _ {0} \in \mathrm{Aut}   \mathrm{G} _ {0} \\ w _ {i} & = x _ {i} ^ {- 1} w (x _ {i}) \in \mathrm{G} _ {0} \quad (1 \leqslant i \leqslant n) \\ \zeta (w) & = ((w _ {1}, \ldots , w _ {n}), w _ {0}) \in \mathrm{P}. \end{array}
\]

对所有 \(w, w'\) 属于 \(\mathrm{Aut}_1\mathrm{G}\)，

\[
\begin{array}{r l} \zeta (w) \zeta (w ^ {\prime}) & = \left((w _ {1}, \ldots , w _ {n}), w _ {0}\right) \left((w _ {1} ^ {\prime}, \ldots , w _ {n} ^ {\prime}), w _ {0} ^ {\prime}\right) \\ & = \left((w _ {1} w _ {0} (w _ {1} ^ {\prime}), \ldots , w _ {n} w _ {0} (w _ {n} ^ {\prime})), w _ {0} w _ {0} ^ {\prime}\right) \\ & = \left((x _ {1} ^ {- 1} w (x _ {1}) w (x _ {1} ^ {- 1} w ^ {\prime} (x _ {1})), \ldots , x _ {n} ^ {- 1} w (x _ {n}) w (x _ {n} ^ {- 1} w ^ {\prime} (x _ {n}))), w _ {0} w _ {0} ^ {\prime}\right) \\ & = \left(((w w ^ {\prime}) _ {1}, \ldots , (w w ^ {\prime}) _ {n}), (w w ^ {\prime}) _ {0}\right) \\ & = \zeta (w w ^ {\prime}) \end{array}
\]

从而 \(\zeta\) 是从 \(\mathrm{Aut}_1\mathrm{G}\) 到 P 的同态。这个同态显然是单射。

我们证明 \(\zeta(\mathrm{Aut}_1\mathrm{G})\) 在 P 中是闭的。令 \(\Phi\) 是 \(\mathrm{Aut}_1\mathrm{G}\) 上的滤子，使得 \(\zeta(\Phi)\) 收敛到 P 中的一点 \(((w_1,\dots,w_n),w_0)\)。于是 \(\Phi\) 逐点收敛到从 G 到 G 的映射 \(v\)。显然，\(v\) 是群 G 的端同态。此外，\(v\) 保持模 \(\mathbf{G}_0\) 的每个陪集稳定，且 \(v |\mathrm{G}_0 = w_0\)。因此 \(v\in \mathrm{Aut}_1\mathrm{G}\)。由于 \(\zeta(v) = ((w_1,\dots,w_n),w_0)\)，我们证明了 \(\zeta(\mathrm{Aut}_1\mathrm{G})\) 在 P 中是闭的。

\begin{enumerate}
\def\labelenumi{(\alph{enumi})}
\setcounter{enumi}{3}
\tightlist
\item
  在证明的 (d) 部分，我们假设 \(\mathbf{K} = \mathbf{R}\)。由第 8 节第 2 号定理 2，\(\zeta(\mathrm{Aut}_1\mathrm{G})\) 是 P 的李子群。将 \(\zeta(\mathrm{Aut}_1\mathrm{G})\) 上的实李群结构传递给 \(\mathrm{Aut}_1\mathrm{G}\)，所用映射为 \(\zeta^{-1}\)。于是 \(\mathrm{Aut}_1\mathrm{G}\) 成为有限维李群。
\end{enumerate}

设 M 是解析流形，f 是从 M 到 \(Aut_{1}G\) 的映射，\(\phi\) 是映射 \((m,g)\mapsto f(m)g\) 从 \(M\times G\) 到 G。有以下等价关系：

\(f\) 是解析的 \(\Leftrightarrow\) 映射 \(m\mapsto (f(m))_i\)，其中 \(0\leqslant i\leqslant n\)，是解析的 \(\Leftrightarrow \left\{ \begin{array}{ll}\text{映射 } m\mapsto f(m)x_i\text{ 从 M 到 G，对 } 1\leqslant i\leqslant n,\text{ 是解析的} \\ \text{且} \\ \text{映射 } (m,g)\mapsto f(m)g\text{ 从 M} \times \mathrm{G}_0\text{ 到 G 是解析的} \end{array} \right.\) \(\Leftrightarrow \phi\) 是解析的。

对于 \(w \in \mathrm{Aut}_1\mathrm{G}\)，有 \(\mathrm{L}(w) = \mathrm{L}(w_0)\)，因此态射 \(w \mapsto \mathrm{L}(w)\) 从 \(\mathrm{Aut}_1\mathrm{G}\) 到 \(\mathrm{Aut}\, \mathrm{L}(\mathrm{G})\) 是解析的。如同 (b) 中一样，\(\mathrm{Aut}_1\mathrm{G}\) 在 \(\mathrm{G}\) 上的作用律所对应的无穷小运算律，是从 \(\mathrm{L}(\mathrm{Aut}_1\mathrm{G})\) 到 \(\mathfrak{a}\) 的同构。

令 C 是 \(G_{0}\) 的一个紧子集，生成 \(G_{0}\)。对于滤子 \(\Phi\)，\(Id_{G}\) 是它在 \(Aut_{1}G\) 上收敛到的极限的充要条件是，\(\Phi|(C \cup \{x_{1}\} \cup \cdots \cup \{x_{n}\})\) 和 \(\Phi^{-1}|(C \cup \{x_{1}\} \cup \cdots \cup \{x_{n}\})\) 一致收敛到

\[
\operatorname{Id} _ {\mathbb {G}} \left| \left(\mathrm{C} \cup \{x _ {1} \} \cup \dots \cup \{x _ {n} \}\right). \right.
\]

因此 \(\mathrm{Aut}_1\mathrm{G}\) 的拓扑就是拓扑 \(\mathcal{T}_{\beta}\)。

显然，\(\mathrm{Aut}_1\mathrm{G}\) 在带有拓扑 \(\mathcal{T}_{\beta}\) 的 Aut G 中是开的。Aut G 上存在与该拓扑相容的李群结构，并在 \(\mathrm{Aut}_1\mathrm{G}\) 上诱导出上述构造的结构（第 8 节第 1 号定理 1 的推论 2）。李群 Aut G 具有定理所述性质这一事实，由 \(\mathrm{Aut}_1\mathrm{G}\) 的相应性质推出。

\begin{enumerate}
\def\labelenumi{(\alph{enumi})}
\setcounter{enumi}{4}
\tightlist
\item
  在证明的 (e) 部分，我们假设 \(\mathbf{K} = \mathbf{C}\)。由 (c) 以及第 8 节第 2 号定理 2，\(\mathrm{Aut}_1\mathrm{G}\) 上存在实李群结构，使得 \(\zeta\) 是从 \(\mathrm{Aut}_1\mathrm{G}\) 到 \(\mathbb{P}\) 的一个实李子群的同构。
\end{enumerate}

\((w, g) \mapsto wg\) 在 \((\mathrm{Aut}_1\mathrm{G}) \times \mathrm{G}\) 上、即在 \(\mathbf{G}\) 上的运算律是实解析的。令 \(\mathbf{D}\) 为相应的无穷小运算律。由第 1 号命题 1 和 2，\(\mathrm{D}(\mathrm{L}(\mathrm{Aut}_1\mathrm{G})) = \mathfrak{a}\)。

对所有 \(\alpha \in \mathfrak{a}\)，令 \(\alpha_0\) 表示 \(\alpha\) 在 \(\mathrm{G}_0\) 上的限制；它是 \(\mathrm{G}_0\) 的无穷小自同构，由于证明的 (b) 部分，我们将它与 \(\mathrm{L}(\mathrm{Aut}\mathrm{G}_0)\) 中的一个元素识别。对 \(1 \leqslant i \leqslant n\)，记

\[
\alpha_ {i} = x _ {i} ^ {- 1} \alpha (x _ {i}) \in \mathrm{L} (\mathrm{G}) = \mathrm{L} (\mathrm{G} _ {0}).
\]

最后，记 \(f(\alpha) = ((\alpha_1, \ldots, \alpha_n), \alpha_0) \in \mathrm{L}(\mathrm{P})\)。则 \(f\) 是一个 \(\mathbf{C}\)-线性映射，作用于 \(\mathfrak{a}\) 并取值于 \(\mathrm{L}(\mathrm{P})\)。

另一方面，显然 \(\mathrm{L}(\zeta) = f\circ \mathrm{D}\)。因此 \(\mathrm{L}(\zeta)(\mathrm{L}(\mathrm{Aut}_1\mathrm{G})) = f(a)\) 是 \(\mathrm{L}(P)\) 的复向量子空间。由第 4 节第 2 号命题 2，\(\zeta (\mathrm{Aut}_1\mathrm{G})\) 是 \(P\) 的复李子群，我们可以完全按照 (d) 中的方式进行：将 \(\zeta (\mathrm{Aut}_1\mathrm{G})\) 上的复李群结构传递给 \(\mathrm{Aut}_1\mathrm{G}\)，所用映射为 \(\zeta^{-1}\)，并如 (d) 中那样看出，\(\mathrm{Aut}_1\mathrm{G}\) 具有与定理的性质 (i)、(ii)、(iii)、(v) 和 (vi) 相应的性质。

显然 \(\mathrm{Aut}_1\mathrm{G}\) 在 Aut G 中带有拓扑 \(\mathcal{T}_{\beta}\) 时是开的。令 \(w\in \mathrm{Aut}G\)。令 \(\sigma\) 为自同构 \(v\mapsto wvw^{-1}\)，作用在 \(\mathrm{Aut}_1\mathrm{G}\) 上。它是实解析的，\(\mathrm{L}(\sigma)\) 是一个 \(\mathbf{R}\)-自同构 \(\mathrm{L}(\mathrm{Aut}_1\mathrm{G})\)，并且

\[
\mathrm{D} \circ \mathrm{L} (\mathrm{Aut} _ {1} \mathrm{G}) \circ \mathrm{D} ^ {- 1}
\]

是 a 的一个 R-自同构。这个自同构也是由 w 通过结构传递得到的 a 的自同构；由于 w 是 K-解析的，我们看到 \(L(\sigma)\) 是 K-线性的。因此 \(\sigma\) 是 K-解析的（第 \(\S\) 3 节第 8 号命题 32）。由第 \(\S\) 1 节第 9 号命题 18，Aut G 上存在唯一的李 K-群结构，使得 \(Aut_1G\) 是 Aut G 的开李子群。这一结构具有定理所述性质这一事实，由 \(Aut_1G\) 的相应性质推出。

推论 1. 设 G 是有限维实李群，\(G_{0}\) 是其单位元分支。假设 G 由 \(G_{0}\) 和有限个元素生成。则 Aut G 具有拓扑 \(T_{\beta}\)，并且是有限维实李群。

推论 2. 设 G 是半单连通实或复李群。群 Int G 是 Aut G 的单位元分支。

映射 \(u \mapsto \mathrm{L}(u)\) 是从 Aut G 到 Aut L(G) 的一个李子群的同构（定理 1）。Int G 在此同构下的像是 Ad G。但 Ad G 是 Aut L(G) 的单位元分支（第 9 节第 8 号命题 30(ii)）。

